\section{Solution methods and verification}\label{sec:methods}
\subsection{Analytical methods}
For linear problems with constant coefficients, the classical analytical tools carry over from elasticity and heat conduction with only the algebra changed by the coupling.
\begin{itemize}
\item \textbf{Plane-wave and normal-mode analysis.} For unbounded or periodic media, \eqref{eq:disp} gives complete frequency-domain information. Real frequency with complex wavenumber describes spatial decay under steady excitation; real wavenumber with complex frequency describes temporal decay of a free disturbance. The two attenuation measures coincide only for weak damping ($Q^{-1}\ll1$), and the relation involves the group velocity: $\Imag\omega\approx-v_g\Imag k$. Using the wrong measure when comparing with an experiment is a common source of apparent inconsistency.
\item \textbf{Integral transforms.} For half-spaces, layers and cylinders, Laplace transformation in time and Fourier (or Hankel) transformation in space reduce the problem to ordinary differential equations with exponential solutions, as in Sections~\ref{sec:transient} and~\ref{sec:surface}. Carrying the heat-flux law as the function $G(s)$ means that one derivation covers the whole template.
\item \textbf{Potential methods and eigenfunction expansion.} Helmholtz potentials decouple shear from dilatation in the bulk, and eigenfunction expansion in the transform domain is often used for finite domains \citep{Singh2020a}.
\item \textbf{Wavefront expansions.} Discontinuities propagate along characteristics determined by the high-frequency limit; \eqref{eq:fronts} is an example. Jump amplitudes follow transport equations along the front and decay exponentially in LS-type theories.
\end{itemize}

\subsection{Numerical inversion of Laplace transforms}\label{sec:inversion}
Closed forms in the transform domain are typically inverted numerically. The commonly used algorithms are the Stehfest method \citep{Stehfest1970}, Durbin's Fourier-series method \citep{Durbin1974} (and the closely related method of Dubner and Abate \citep{DubnerAbate1968}), the Honig--Hirdes method \citep{HonigHirdes1984}, and contour-deformation methods such as Talbot's \citep{Talbot1979}; Abate and Whitt \citep{AbateWhitt2006} review the common framework. Because applications to thermoelastic problems often report numbers from one of these methods without a check, we benchmarked four implementations (Talbot, Stehfest, de Hoog et al.\ \citep{deHoog1982}, and Cohen acceleration \citep{Cohen2007}, as implemented in a standard multiprecision library with default parameters at 15 digits) on four transforms with known inverses, chosen to represent the features of thermoelastic responses: a smooth decay, an oscillation, a discontinuous step with a delay, and a diffusive response (\Cref{tab:bench}).
\begin{table}[ht]\centering\small
\caption{Maximum absolute error of four numerical Laplace inversion methods on test transforms (evaluated at $t=0.5,2,5,10$; for the step at $t=0.5,0.95,1.05,2,5$). Default parameters, 15 digits. ``$4\times10^{193}$'' indicates a numerical blow-up.}\label{tab:bench}
\begin{tabular}{@{}lrrrr@{}}\toprule
Test function & Talbot & Stehfest & de Hoog & Cohen\\\midrule
smooth: $1/(s+1)^2$ & 6.9e$-$18 & 7.1e$-$13 & 9.8e$-$18 & 6.9e$-$18\\
oscillatory: $1/\sqrt{s^2+1}$ & 3.9e$-$2 & 6.6e$-$4 & 3.6e$-$16 & 6.1e$-$16\\
step: $e^{-s}/s$ & 3.8e193 & 1.8e$-$1 & 1.4e$-$6 & 1.8e$-$1\\
diffusive: $e^{-\sqrt s}/s$ & 1.1e$-$16 & 5.5e$-$14 & 1.1e$-$16 & 1.1e$-$16\\
\bottomrule\end{tabular}
\end{table}
Four points are worth recording. (i) For smooth and diffusive responses all four methods are essentially exact ($<10^{-12}$), which is why most published diffusion-dominated results are insensitive to the method. (ii) For the step with a delay, Talbot's method \emph{fails catastrophically} at times before the delay ($t=0.95<1$): its contour, which extends into the left half-plane, is multiplied by $e^{-s}$ and grows without bound. This is exactly the situation of a Cattaneo front, whose transform contains $e^{-x/c_T\cdot s}$. The remedy used in Section~\ref{sec:transient} is to set the solution to zero ahead of the known front and invert only for $t>x/c_T$. (iii) Stehfest and Cohen give errors of order $0.2$ near the discontinuity (they smear the front and ring), while the de Hoog method yields $1.4\times10^{-6}$ at the tested points. (iv) For oscillatory functions the default Talbot setting is inadequate at $t=10$ (error $0.04$), while de Hoog and Cohen are accurate. The methods can be tuned (number of terms, contour parameter), and the table is not a ranking; its message is that the choice and the checks matter. We recommend: invert with two independent methods, compare against a known limit, and treat results near discontinuities as unreliable beyond the known front position.

\subsection{Time-domain numerics}
For problems that do not admit transform solutions (heterogeneous, nonlinear, complex geometry) finite-element and finite-difference methods in space and time are used. Two issues are specific to hyperbolic thermal models. First, schemes designed for parabolic heat equations (implicit, with numerical dissipation) smear the thermal front and then cannot represent the finite-speed effect that motivates the model; explicit or high-resolution schemes with controlled dissipation, or space--time discontinuous Galerkin formulations, are appropriate. Second, the time step must resolve the relaxation time ($\Delta t\ll\tau_q$) and satisfy a CFL condition based on the \emph{largest} wave speed, which is $\max(c_1,c_T)$ and not $c_1$ alone: when $c_T\gg c_1$ (as in many metals) the thermal front dominates the stability bound. Memory-dependent and fractional operators require storage of history and, in an unoptimised implementation, an $O(N^2)$ cost in the number of time steps; short-memory approximations and fast convolution reduce this, and should be accompanied by an error estimate. Finite-element implementations with memory-dependent heat transfer in fibre-reinforced cylinders \citep{Sur2019} and rotating nano-beams \citep{Sur2025} illustrate the approach.

\subsection{A verification discipline}
We find that the following checks catch most errors in new thermoelastic wave calculations, and we recommend that papers state which they performed.
\begin{enumerate}
\item \textbf{Reduce to a known limit}: set the coupling to zero (elastic and thermal problems separately), and set the lags to zero (Fourier).
\item \textbf{Check the dispersion relation symbolically} against the field equations, as in Appendix~\ref{app:disp}; sign and factor errors in the coupling terms are detected immediately.
\item \textbf{Check energy}: for conservative members of the family (GN-II) the energy is constant; for the others, the energy balance \eqref{eq:energy} must hold with the stated dissipation; for interfaces, the net flux must be continuous (Section~\ref{sec:interfaces}).
\item \textbf{Check high- and low-frequency asymptotics}: adiabatic and isothermal speeds, front speeds (e.g.\ \eqref{eq:fronts}), classical diffusive scaling.
\item \textbf{Check branches}: track roots by continuity and ensure that Taylor truncation has not introduced spurious roots.
\item \textbf{Report dimensionless groups}: $\varepsilon$, $\omega/\omega_*$, $\omega\tau$, $c_T/c_1$, $\xi\omega/c_1$, so that results can be compared across studies.
\end{enumerate}


\begin{summarybox}[Summary]
Solution methods: analytical separation for simple geometries, Laplace/Fourier inversion for transients, numerical root finding for dispersion and secular equations. Verification discipline: reduce to the elastic limit, check Snell invariants, verify net-flux continuity (not mode-flux sums), test boundary conditions by substitution, state the flux definition for lagged and memory models. Report dimensionless groups $\varepsilon$, $\omega/\omega_*$, $\omega\tau$, $c_T/c_1$, $\xi\omega/c_1$ for comparability.
\end{summarybox}